\documentclass{article}
\usepackage{algorithm, algorithmic, amsfonts, amsmath}

\begin{document}
\begin{algorithm}[!h]
\renewcommand{\algorithmicrequire}{\textbf{Input:}}
\renewcommand{\algorithmicensure}{\textbf{Output:}}
	\caption{IW$\ell$ (importance reweighting for classification with noisy labels)}
	\begin{algorithmic}[1]
	    \REQUIRE
	    Noisy sample $\{(X_{i},\hat{Y}_{i})\}_{i=1}^{n}\sim D_{\rho}$;
	    any surrogate loss $\ell$ designed for ordinary classification;
	    hypothesis class $\mathcal{F}$
	    \ENSURE
	    Classifier $\hat{f}_{n}$ consistent for the clean distribution $D$
	    \STATE \textbf{Setting:} asymmetric random classification noise --- a true $+1$ flips with
	    rate $\rho_{+1}$ and a true $-1$ with rate $\rho_{-1}$; PU is the special case
	    $\rho_{-1}=0$ % PU is one-sided noise, so the method applies to PU directly
	    \STATE \textbf{Step 1: estimate the noisy posterior and the noise rates}
	    \STATE Fit $P_{D_{\rho}}(\hat{Y}\mid X)$ on the noisy sample % estimable from the noisy data alone
	    \STATE Estimate $\rho_{+1},\rho_{-1}$ (the paper bounds them by the conditional
	    probability and estimates them from the fitted posterior)
	    \STATE \textbf{Step 2: importance weights}
	    \STATE $\beta(X,\hat{Y})=\dfrac{P_{D}(Y\mid X)}{P_{D_{\rho}}(\hat{Y}\mid X)}
	    =\dfrac{P_{D_{\rho}}(\hat{Y}\mid X)-\rho_{-\hat{Y}}}
	    {\left(1-\rho_{+1}-\rho_{-1}\right)P_{D_{\rho}}(\hat{Y}\mid X)}$ % the density ratio between clean and noisy distributions, expressible in closed form from the noise rates
	    \STATE Set $\beta(X,\hat{Y})=0$ whenever $P_{D_{\rho}}(\hat{Y}\mid X)=0$ % a convention that keeps the weight non-negative
	    \STATE \textbf{Step 3: minimise the reweighted empirical risk}
	    \STATE $\hat{f}_{n}=\arg\min_{f\in\mathcal{F}}\ \dfrac{1}{n}\sum_{i=1}^{n}
	    \beta(X_{i},\hat{Y}_{i})\,\ell\!\left(f(X_{i}),\hat{Y}_{i}\right)$ % any off-the-shelf loss carries over unchanged; only the weighting is added
		\RETURN $\hat{f}_{n}$
	\end{algorithmic}
\end{algorithm}
% Liu \& Tao, IEEE TPAMI 2016, "Classification with Noisy Labels by Importance
% Reweighting". Compiled from equations (1) to (3), Lemma 2 and Theorem 1 of the original.
% Position: one of the key theoretical works placing PU learning inside the framework of one-sided label noise.
% Central result (Theorem 1): any surrogate loss designed for ordinary classification can be used under asymmetric RCN,
% with the same consistency rate as the corresponding ordinary algorithm, provided it is weighted by beta and provided
% the noisy posterior P_{D_rho}(Y_hat|X) and the noise rates rho_{+1}, rho_{-1} are estimated accurately enough.
% The cost sits in that proviso: the method trades the problem of designing a new loss for the problem of estimating
% two noise rates. The form of beta shows why, since its denominator carries (1 - rho_+1 - rho_-1) and the weight
% explodes as the two rates sum toward 1. The survey places it in branch H, since its direct product is a noise-rate estimate.
\end{document}
